\section{Σύνοψη Θεωρητικού Υποβάθρου}

\subsection{Εισαγωγή}

Το κεφάλαιο αυτό συνοψίζει τις βασικές έννοιες και τεχνολογίες που παρουσιάστηκαν στο θεωρητικό υπόβαθρο της παρούσας εργασίας. Στόχος είναι η παροχή μιας συνεκτικής επισκόπησης των θεμελίων που υποστηρίζουν την προτεινόμενη προσέγγιση.

\subsection{Πολυπλοκότητα Δικτύων \en{5G} και \en{6G}}

Τα σύγχρονα δίκτυα επικοινωνιών έχουν εξελιχθεί σημαντικά, με τα συστήματα \en{5G} και τα μελλοντικά 6G να εισάγουν αυξημένη πολυπλοκότητα, υψηλότερους ρυθμούς δεδομένων και μαζική συνδεσιμότητα. Τα χαρακτηριστικά αυτά δημιουργούν νέες προκλήσεις στη διαχείριση και την ασφάλεια των δικτύων.

\subsection{Σημασία Αυτοματοποιημένης Ανάλυσης Αρχείων Καταγραφής}

Τα αρχεία καταγραφής δικτύου αποτελούν κρίσιμη πηγή πληροφορίας για την παρακολούθηση και την ασφάλεια των συστημάτων. Ωστόσο, ο μεγάλος όγκος και η μη δομημένη φύση τους καθιστούν τη χειροκίνητη ανάλυση μη πρακτική. Συνεπώς, η αυτοματοποιημένη ανάλυση είναι απαραίτητη για την ανίχνευση ανωμαλιών και την άμεση απόκριση σε κυβερνοαπειλές.

\subsection{Ρόλος των \en{LLMs} στην Κυβερνοασφάλεια}

Τα Μεγάλα Γλωσσικά Μοντέλα (\en{LLMs}) έχουν επιδείξει ισχυρές δυνατότητες στην επεξεργασία μη δομημένων δεδομένων, καθιστώντας τα κατάλληλα για ανάλυση αρχείων καταγραφής και εφαρμογές κυβερνοασφάλειας. Η ικανότητά τους να κατανοούν συμφραζόμενα και να εντοπίζουν σύνθετα πρότυπα επιτρέπει αποτελεσματικότερη ανίχνευση απειλών σε σύγκριση με παραδοσιακές μεθόδους.

\subsection{Ενσωμάτωση Προηγμένων Τεχνικών}

Ο συνδυασμός \en{LLMs} με τεχνικές αποδοτικής προσαρμογής παραμέτρων όπως η \en{LoRA} και κατανεμημένα πλαίσια μάθησης όπως η \en{Federated Learning} παρέχει μια επεκτάσιμη και φιλική προς την ιδιωτικότητα λύση. Η ενοποίηση αυτή αντιμετωπίζει τους περιορισμούς των επιμέρους προσεγγίσεων και επιτρέπει αποδοτική εφαρμογή σε σύγχρονα δίκτυα.

\subsection{Συμπέρασμα}

Το θεωρητικό υπόβαθρο που παρουσιάστηκε αναδεικνύει τη σημασία του συνδυασμού προηγμένων τεχνικών μηχανικής μάθησης με κατανεμημένα συστήματα για την αντιμετώπιση των προκλήσεων των δικτύων επόμενης γενιάς. Οι έννοιες αυτές αποτελούν τη βάση για τη μεθοδολογία και την πειραματική αξιολόγηση που ακολουθεί.